\subsubsection{Toeplitz 负惯性的最小采样阈值、正定铅笔重建与一素数审计}\label{subsubsec:xi-toeplitz-minimal-sampling-gramshift-one-prime}
在 Toeplitz--扫描 Hankel 同构分解（推论 \ref{cor:xi-toeplitz-scan-hankel-congruence-decomposition}）的稳定区间内，Toeplitz 层级的负方向由 Hankel 指纹块 \(\mathsf H_{\kappa-1}\) 的 Gram \(\mathsf H_{\kappa-1}^\ast\mathsf H_{\kappa-1}\) 完全确定。
本段抽取三类可核验后果：其一给出负惯性判定所需样本长度的严格最小阈值；其二给出完全正定的 Gram--shift 广义特征值铅笔以重建缺陷节点谱；其三给出将满秩核验压缩为单个模素数证书的算术审计接口，并把连续缺陷熵与离散负块谱质量建立不可压缩的定量链接。

\begin{theorem}[负惯性判定的最小采样阈值]\label{thm:xi-toeplitz-negative-inertia-minimal-sampling}
在有限缺陷口径下，令扫描谱指纹为 \(\{u_n\}_{n\ge 0}\) 并定义 Hankel 窗口
\[
\mathsf H_M:=(u_{p+q})_{0\le p,q\le M}.
\]
假设进入稳定区间：存在 \(N_0\) 使对一切充分大 \(N\) 有 \(\nu_-(\mathbf{T}_N(C))=\kappa\)，并且 \(\mathrm{rank}(\mathsf H_{\kappa-1})=\kappa\)（从而推论 \ref{cor:xi-toeplitz-scan-hankel-congruence-decomposition} 可用）。
则：
\begin{enumerate}
  \item Hankel 指纹块 \(\mathsf H_{\kappa-1}\) 的全部条目仅依赖有限样本族 \(\{u_n\}_{0\le n\le 2\kappa-2}\)；
  因而 Toeplitz 层级的负惯性与负子空间在稳定区间内可由该有限样本族完全确定。
  \item 该样本长度在一般位置意义下严格最小：存在 Zariski 开稠密集合，使得任意（允许自适应、允许随机化的）判定器若仅观测 \(\{u_n\}_{0\le n\le 2\kappa-3}\)，则在等权二点对抗先验下的最小错误概率至少为 \(1/2\)；从而不能以该前缀普适决定 \(\mathrm{rank}(\mathsf H_{\kappa-1})\)（亦不能决定 \(\nu_-(\mathbf{T}_N(C))\) 是否等于 \(\kappa\)）。
\end{enumerate}
\end{theorem}
\begin{proof}
第一条由定义直接给出：\((\mathsf H_{\kappa-1})_{p+1,q+1}=u_{p+q}\)，其中 \(0\le p+q\le 2\kappa-2\)。
在稳定区间内，推论 \ref{cor:xi-toeplitz-scan-hankel-congruence-decomposition} 给出
\(
\nu_-(\mathbf{T}_N(C))=\mathrm{rank}(\mathsf H_{\kappa-1})
\)
以及负子空间的合同刻画，因此负惯性与负子空间均由 \(\mathsf H_{\kappa-1}\) 决定，从而由 \(\{u_n\}_{0\le n\le 2\kappa-2}\) 决定。
\par
对第二条，固定任意前缀数据 \((u_0,\dots,u_{2\kappa-3})\)。
注意 \(\det(\mathsf H_{\kappa-1})\) 作为变量 \(u_{2\kappa-2}\) 的函数是一次仿射的：\(u_{2\kappa-2}\) 仅出现在右下角条目 \((\kappa,\kappa)\)。
因此存在仅依赖 \((u_0,\dots,u_{2\kappa-3})\) 的常数 \(A,B\) 使
\[
\det(\mathsf H_{\kappa-1})=A\,u_{2\kappa-2}+B.
\]
当 \(A\neq 0\) 时，既可选择 \(u_{2\kappa-2}\) 使 \(\det(\mathsf H_{\kappa-1})=0\)，亦可选择另一 \(u'_{2\kappa-2}\) 使 \(\det(\mathsf H_{\kappa-1})\neq 0\)；
而两种选择在 \(\{u_n\}_{0\le n\le 2\kappa-3}\) 上完全一致。
系数 \(A\) 是相应余子式的多项式函数，故 \(\{A\neq 0\}\) 为 Zariski 开集；只要存在一个实例使 \(A\neq 0\)，该开集即非空从而开稠密。
于是一般位置下，\(\{u_n\}_{0\le n\le 2\kappa-3}\) 无法决定 \(\det(\mathsf H_{\kappa-1})\) 的消失/非消失，进而无法决定 \(\mathrm{rank}(\mathsf H_{\kappa-1})\)；并且对任何仅依赖该前缀的判定器（含随机化），上述两种末端延拓诱导完全相同的观测分布，故在等权二点先验下其最小错误概率至少为 \(1/2\)。由推论 \ref{cor:xi-toeplitz-scan-hankel-congruence-decomposition}，同一不可判定性等价传递到 Toeplitz 负惯性是否达到 \(\kappa\)。证毕。
\end{proof}
\leanverified{paper\_xi\_toeplitz\_negative\_inertia\_minimal\_sampling}

\begin{proposition}[Hankel 最小采样的“一点例外”结构与鲁棒判定]\label{prop:xi-hankel-minimal-sampling-one-point-exception-robust}
\leanverified{paper\_xi\_hankel\_minimal\_sampling\_one\_point\_exception\_robust}
沿用定理 \ref{thm:xi-toeplitz-negative-inertia-minimal-sampling} 的记号，并额外假设
\(\det(\mathsf H_{\kappa-2})\neq 0\)。
则存在唯一临界值 \(u_{2\kappa-2}^{\star}\in\CC\)（仅依赖于 \(\{u_n\}_{0\le n\le 2\kappa-3}\)），使得
\[
\det(\mathsf H_{\kappa-1})
=\det(\mathsf H_{\kappa-2})\,\bigl(u_{2\kappa-2}-u_{2\kappa-2}^{\star}\bigr).
\]
因此
\[
\mathrm{rank}(\mathsf H_{\kappa-1})=\kappa
\quad\Longleftrightarrow\quad
u_{2\kappa-2}\neq u_{2\kappa-2}^{\star}.
\]
并且该判定具有定量裕度：若 \(|u_{2\kappa-2}-u_{2\kappa-2}^{\star}|\ge \eta\)，则
\[
\bigl|\det(\mathsf H_{\kappa-1})\bigr|\ \ge\ \bigl|\det(\mathsf H_{\kappa-2})\bigr|\,\eta.
\]
\end{proposition}
\begin{proof}
沿最后一行最后一列按 Laplace 展开，\(\det(\mathsf H_{\kappa-1})\) 关于变量 \(u_{2\kappa-2}\) 为仿射函数，且其系数等于删去最后一行最后一列所得主子式行列式 \(\det(\mathsf H_{\kappa-2})\)。
故存在仅依赖 \(\{u_n\}_{0\le n\le 2\kappa-3}\) 的常数 \(B\) 使
\[
\det(\mathsf H_{\kappa-1})=\det(\mathsf H_{\kappa-2})\,u_{2\kappa-2}+B.
\]
在 \(\det(\mathsf H_{\kappa-2})\neq 0\) 下，令
\(
u_{2\kappa-2}^{\star}:=-B/\det(\mathsf H_{\kappa-2})
\)
即得第一式与唯一性。
秩判别由
\(
\det(\mathsf H_{\kappa-1})\neq 0\iff \mathrm{rank}(\mathsf H_{\kappa-1})=\kappa
\)
给出；裕度不等式由第一式直接推出。证毕。
\end{proof}

\begin{theorem}[Gram--shift 正定铅笔的缺陷节点谱]\label{thm:xi-gram-shift-positive-pencil-node-spectrum}
在定理 \ref{thm:xi-toeplitz-negative-inertia-minimal-sampling} 的稳定区间设定下，定义两块移位 Hankel 窗口
\[
H_0:=(u_{p+q})_{0\le p,q\le \kappa-1}=\mathsf H_{\kappa-1},
\qquad
H_1:=(u_{p+q+1})_{0\le p,q\le \kappa-1}.
\]
并令
\[
G_0:=H_0^\ast H_0\succ0,\qquad G_1:=H_0^\ast H_1.
\]
定义移位实现矩阵 \(A:=H_0^{-1}H_1\)。
则：
\begin{enumerate}
  \item 存在严格恒等式 \(G_1=G_0A\)，从而 \(A=G_0^{-1}G_1\)；
  并且 \(A\) 的谱等于正定铅笔 \((G_1,G_0)\) 的广义特征值集合。
  \item 若进一步满足有限指数分解 \eqref{eq:xi-scan-fingerprint-exponential-sum}（定理 \ref{thm:xi-dethankel-vandermonde-square-finite-blaschke} 的设定），则 \(A\) 与对角阵 \(\mathrm{diag}(z_1,\dots,z_\kappa)\) 相似，且 \(\sigma(A)=\{z_1,\dots,z_\kappa\}\) 给出缺陷节点谱。
  因而节点重建可归约为求解完全正定的广义特征值问题 \(G_1v=\lambda G_0v\)。
\end{enumerate}
\end{theorem}
\begin{proof}
由 \(H_1=H_0A\) 左乘 \(H_0^\ast\) 得 \(G_1=H_0^\ast H_1=H_0^\ast H_0A=G_0A\)，从而 \(A=G_0^{-1}G_1\)。
若 \(Av=\lambda v\)，则 \(G_1v=G_0Av=\lambda G_0v\)，故 \(\lambda\) 为铅笔 \((G_1,G_0)\) 的广义特征值；反向同理。
\par
在 \eqref{eq:xi-scan-fingerprint-exponential-sum} 下，令 Vandermonde 矩阵 \(V\) 如定理 \ref{thm:xi-dethankel-vandermonde-square-finite-blaschke}，并记 \(W:=\mathrm{diag}(w_1,\dots,w_\kappa)\)、\(Z:=\mathrm{diag}(z_1,\dots,z_\kappa)\)。
直接计算得
\[
H_0=(4\pi)\,VWV^{\mathsf T},
\qquad
H_1=(4\pi)\,VWZV^{\mathsf T}.
\]
从而
\[
A=H_0^{-1}H_1
=\bigl(V^{-{\mathsf T}}W^{-1}V^{-1}\bigr)\,\bigl(VWZV^{\mathsf T}\bigr)
=V^{-{\mathsf T}}ZV^{\mathsf T},
\]
故 \(A\) 与 \(Z\) 相似并具有相同谱 \(\{z_j\}\)。证毕。
\end{proof}

\begin{theorem}[Hankel 行列式的 Gram--shift 判别式恒等式]\label{thm:xi-hankel-determinant-gramshift-discriminant-identity}
\leanverified{paper\_xi\_hankel\_determinant\_gramshift\_discriminant\_identity}
在定理 \ref{thm:xi-gram-shift-positive-pencil-node-spectrum} 的稳定区间设定下，
若进一步满足有限指数分解 \eqref{eq:xi-scan-fingerprint-exponential-sum}，
则对 \(A:=H_0^{-1}H_1\) 的特征多项式
\[
\chi_A(\lambda):=\det(\lambda I_\kappa-A)=\prod_{j=1}^{\kappa}(\lambda-z_j)
\]
的（无符号）判别式
\[
\operatorname{Disc}(\chi_A):=\prod_{1\le i<j\le \kappa}(z_j-z_i)^2
\]
成立严格恒等式
\[
\det(H_0)
=(4\pi)^{\kappa}\Bigl(\prod_{j=1}^{\kappa} w_j\Bigr)\operatorname{Disc}(\chi_A).
\]
\end{theorem}
\begin{proof}
由定理 \ref{thm:xi-gram-shift-positive-pencil-node-spectrum} 的证明，存在 Vandermonde 矩阵 \(V\) 使
\(
H_0=(4\pi)\,V W V^{\mathsf T}
\)，其中 \(W=\mathrm{diag}(w_1,\dots,w_\kappa)\)。
取行列式得
\[
\det(H_0)=(4\pi)^\kappa\det(V)^2\prod_{j=1}^{\kappa}w_j.
\]
而 \(\det(V)=\prod_{1\le i<j\le\kappa}(z_j-z_i)\)，故 \(\det(V)^2=\operatorname{Disc}(\chi_A)\)，代回即得结论。证毕。
\end{proof}

\begin{proposition}[Vandermonde 对角化的不可规避病态下界]\label{prop:xi-vandermonde-conditioning-lower-bound-disc}
\leanverified{paper\_xi\_vandermonde\_conditioning\_lower\_bound\_disc}
在定理 \ref{thm:xi-gram-shift-positive-pencil-node-spectrum} 的设定下，
设 \(V\in\CC^{\kappa\times\kappa}\) 为节点 \(\{z_j\}_{j=1}^{\kappa}\) 的 Vandermonde 矩阵
\(
V_{ij}:=z_j^{\,i-1}
\)（\(1\le i,j\le \kappa\)），并令 \(r:=\max_j|z_j|\)。
则其 \(2\)-范数条件数满足
\[
\kappa_2(V)
\ge \frac{\max_j\|{\rm col}_j(V)\|_2}{|\det(V)|^{1/\kappa}}
\ge \frac{\sqrt{\sum_{i=0}^{\kappa-1}r^{2i}}}{|\operatorname{Disc}(\chi_A)|^{1/(2\kappa)}}.
\]
\end{proposition}
\begin{proof}
记 \(V\) 的奇异值为 \(\sigma_1\ge\cdots\ge\sigma_\kappa>0\)。
由 \(|\det(V)|=\prod_{i=1}^{\kappa}\sigma_i\) 与 \(\sigma_i\ge \sigma_\kappa\) 得
\(
\sigma_\kappa^\kappa\le |\det(V)|
\)，从而 \(\sigma_\kappa\le |\det(V)|^{1/\kappa}\)。
又由算子范数定义有 \(\sigma_1=\|V\|_2\ge \max_j\|{\rm col}_j(V)\|_2\)。
故
\[
\kappa_2(V)=\frac{\sigma_1}{\sigma_\kappa}\ge \frac{\max_j\|{\rm col}_j(V)\|_2}{|\det(V)|^{1/\kappa}}.
\]
取 \(|z_j|=r\) 的列可得 \(\|{\rm col}_j(V)\|_2^2=\sum_{i=0}^{\kappa-1}r^{2i}\)。
再由 \(|\det(V)|^2=|\operatorname{Disc}(\chi_A)|\) 得到所示不等式。证毕。
\end{proof}

\begin{theorem}[缺陷节点谱的最短采样窗：\texorpdfstring{$2\kappa$}{2kappa} 的不可约性]\label{thm:xi-node-spectrum-minimal-window-2kappa}
在定理 \ref{thm:xi-gram-shift-positive-pencil-node-spectrum} 的稳定区间设定与有限指数分解 \eqref{eq:xi-scan-fingerprint-exponential-sum}（权重非零、节点两两不同）的情形下，
节点谱 \(\{z_j\}_{j=1}^{\kappa}\) 的重建满足如下最小窗长律：
\begin{enumerate}
  \item \textnormal{（充分性）}长度为 \(2\kappa\) 的采样窗 \((u_0,\dots,u_{2\kappa-1})\) 可构造性确定节点多重集 \(\{z_j\}\)，并且
  \(
  \sigma(A)=\{z_1,\dots,z_\kappa\}
  \)
  其中 \(A=H_0^{-1}H_1\)。
  \item \textnormal{（一般位置下的不可约性）}在 Zariski 开稠密意义下，任何严格短于 \(2\kappa\) 的采样窗都不能唯一确定节点多重集；
  具体地，仅给定 \((u_0,\dots,u_{2\kappa-2})\) 时，对应的参数纤维在一般位置上具有正维度，因而存在无穷多组不同的节点--权重数据给出相同的截断序列。
\end{enumerate}
\end{theorem}
\leanverified{paper\_xi\_node\_spectrum\_minimal\_window\_2kappa}
\begin{proof}
\textbf{(1)} 由定理 \ref{thm:xi-gram-shift-positive-pencil-node-spectrum}，构造 \(H_0\) 仅需 \(\{u_n\}_{0\le n\le 2\kappa-2}\)，构造 \(H_1\) 额外需要 \(u_{2\kappa-1}\)；
当 \(H_0\) 可逆时 \(A=H_0^{-1}H_1\) 可定义且 \(\sigma(A)=\{z_j\}\)。
\par
\textbf{(2)} 令 \(s_n:=u_n/(4\pi)\) 并考虑 Prony 矩映射
\[
\Phi_{\kappa}:\ (z_1,\dots,z_\kappa,w_1,\dots,w_\kappa)\ \longmapsto\ (s_0,\dots,s_{2\kappa-1}),
\qquad
s_n=\sum_{j=1}^{\kappa}w_j z_j^{n}.
\]
由定理 \ref{thm:xi-prony-moment-map-jacobian-delta4}（取 \(k=\CC\)）的 Jacobian 闭式，
在一般位置（\(w_j\neq 0\) 且 \(z_i\neq z_j\)）上有 \(\det D\Phi_\kappa\neq 0\)，从而 \(\Phi_\kappa\) 在该处局部双全纯。
令 \(\pi:\CC^{2\kappa}\to\CC^{2\kappa-1}\) 为遗忘最后一坐标的投影
\(
\pi(s_0,\dots,s_{2\kappa-1})=(s_0,\dots,s_{2\kappa-2})
\)。
则在上述一般位置点处，
\(
\mathrm{rank}\,D(\pi\circ\Phi_\kappa)=2\kappa-1
\)，
由秩定理知 \((\pi\circ\Phi_\kappa)^{-1}(s_0,\dots,s_{2\kappa-2})\) 的纤维在一般位置上为正维（至少为一维）。
\par
最后，若节点 \((z_1,\dots,z_\kappa)\) 固定且两两不同，则权重向量 \((w_1,\dots,w_\kappa)\) 由前 \(\kappa\) 个矩 \((s_0,\dots,s_{\kappa-1})\) 通过 Vandermonde 线性系统唯一确定；
因此上述正维纤维不能仅由权重自由度贡献，节点本身必随之变化，从而节点多重集不能由长度 \(2\kappa-1\) 的采样窗唯一确定。证毕。
\end{proof}

\begin{corollary}[Prony 判定--重建的恰差一采样律]\label{cor:xi-prony-decision-reconstruction-one-sample-gap}
\leanverified{paper\_xi\_prony\_decision\_reconstruction\_one\_sample\_gap}
在定理 \ref{thm:xi-toeplitz-negative-inertia-minimal-sampling} 与定理 \ref{thm:xi-node-spectrum-minimal-window-2kappa} 的设定下，
假设 \((u_n)\) 在一般位置上具有最小 Prony 形式
\(
u_n=\sum_{j=1}^{\kappa}w_j z_j^{n}
\)
（\(w_j\neq 0\)、\(z_i\neq z_j\)）。
则三条结论同时成立：
\begin{enumerate}
  \item \textnormal{（判定阈值）}给定前缀 \((u_0,\dots,u_{2\kappa-2})\)（长度 \(2\kappa-1\)）即可在一般位置下判定 \(\mathrm{rank}(\mathsf H_{\kappa-1})=\kappa\)（等价于 \(\det(\mathsf H_{\kappa-1})\neq 0\)）；而任何仅依赖 \((u_0,\dots,u_{2\kappa-3})\) 的规则在一般位置下不能完成该判定。
  \item \textnormal{（纤维维数）}截断映射
  \[
  (u_n)_{n\ge 0}\longmapsto (u_0,\dots,u_{2\kappa-2})
  \]
  在 Hankel 秩为 \(\kappa\) 的一般位置上具有一维（复）纤维；因此“判定窗”仍保留恰一维自由度。
  \item \textnormal{（重建阈值）}加入单个额外采样 \(u_{2\kappa-1}\)（长度 \(2\kappa\)）即可在一般位置下唯一重建节点多重集 \(\{z_j\}\) 及权重 \(\{w_j\}\)（至多差一个排列）。
\end{enumerate}
从而在一般位置下，“判定”与“重建”之间存在恰差一项采样的不可约间隙。
\end{corollary}

\begin{proof}
第一条由定理 \ref{thm:xi-toeplitz-negative-inertia-minimal-sampling} 直接给出。
第二条由 \(\mathrm{rank}\,D(\pi\circ\Phi_\kappa)=2\kappa-1\)（定理 \ref{thm:xi-node-spectrum-minimal-window-2kappa} 的证明）得到：由于参数空间维数为 \(2\kappa\)，故一般位置纤维维数恰为 \(1\)。
第三条由定理 \ref{thm:xi-node-spectrum-minimal-window-2kappa} 的充分性部分（构造 \(H_0,H_1\) 从而得到 \(A=H_0^{-1}H_1\) 与 \(\sigma(A)=\{z_j\}\)，再线性解出权重）给出。证毕。
\end{proof}

\begin{corollary}[有限窗口的有理重构：节点多项式与谱不变量]\label{cor:xi-node-polynomial-rational-reconstruction}
\leanverified{paper\_xi\_node\_polynomial\_rational\_reconstruction}
沿用定理 \ref{thm:xi-gram-shift-positive-pencil-node-spectrum} 的记号，并假设 \(H_0\) 在 \(\CC\) 上可逆。
令
\[
\Pi_\kappa(T):=\det(TI_\kappa-A)\in\CC[T],
\qquad A:=H_0^{-1}H_1.
\]
则 \(\Pi_\kappa\) 的系数属于有理函数域 \(\QQ(u_0,\dots,u_{2\kappa-1})\)，并且其根集（按代数重数计）恰为节点谱 \(\{z_1,\dots,z_\kappa\}\)。
因此，节点谱的任意对称多项式不变量（例如 \(\sum_j z_j\)、\(\prod_j z_j\) 及 Vandermonde 判别式）均可由有限窗口 \(\{u_n\}_{0\le n\le 2\kappa-1}\) 在纯有理运算内重构。
\end{corollary}
\begin{proof}
由定义，
\(
H_0=(u_{p+q})_{0\le p,q\le \kappa-1}
\)
与
\(
H_1=(u_{p+q+1})_{0\le p,q\le \kappa-1}
\)
的全部条目都属于 \(\{u_n\}_{0\le n\le 2\kappa-1}\)。
当 \(H_0\) 可逆时，\(A=H_0^{-1}H_1\) 的各条目属于 \(\QQ(u_0,\dots,u_{2\kappa-1})\)，从而 \(\Pi_\kappa\) 的系数亦属于该域。
根集结论由定理 \ref{thm:xi-gram-shift-positive-pencil-node-spectrum} 的第二条直接给出。证毕。
\end{proof}

\begin{proposition}[节点多项式的良约化判据：由行列式与判别式封口]\label{prop:xi-node-polynomial-good-reduction-discriminant}
\leanverified{paper\_xi\_node\_polynomial\_good\_reduction\_discriminant}
沿用推论 \ref{cor:xi-node-polynomial-rational-reconstruction} 的记号。
取公共分母 \(D\in\NN\) 使得
\(
\widehat H_0:=DH_0,\ \widehat H_1:=DH_1\in \Mat_\kappa(\ZZ)
\)，并令 \(\widehat\Pi_\kappa\in\ZZ[T]\) 为 \(\Pi_\kappa\) 的一个原始整系数模型（清分母并去除公因子后得到的首一多项式）。
则对任意素数 \(p\nmid D\)，若同时满足
\[
p\nmid \det(\widehat H_0)
\qquad\text{且}\qquad
p\nmid \operatorname{Disc}(\widehat\Pi_\kappa),
\]
则 \(\widehat\Pi_\kappa\bmod p\) 在 \(\FF_p\) 上无重根（等价地为可分多项式），从而节点谱在模 \(p\) 约化下保持分离。
特别地，导致“\(H_0\) 不可逆”或“节点合并”的坏素数集合包含于
\(
\{p:\ p\mid D\det(\widehat H_0)\operatorname{Disc}(\widehat\Pi_\kappa)\}
\)，因而为有限集。
\end{proposition}
\begin{proof}
若 \(p\nmid D\) 且 \(p\nmid \det(\widehat H_0)\)，则 \(H_0\bmod p\) 可逆，从而移位实现矩阵 \(A:=H_0^{-1}H_1\) 在 \(\FF_p\) 上可定义。
另一方面，\(\widehat\Pi_\kappa\bmod p\) 无重根当且仅当
\(
\gcd(\widehat\Pi_\kappa,\widehat\Pi_\kappa')=1
\)
在 \(\FF_p[T]\) 中成立，这等价于判别式 \(\operatorname{Disc}(\widehat\Pi_\kappa)\) 在 \(\FF_p\) 中非零。
故在所述条件下 \(\widehat\Pi_\kappa\bmod p\) 可分，节点约化不发生合并。
最后一条由“坏素数必须整除某个非零整数”得到：若 \(p\) 不整除该积，则同时避开两类坏情形。证毕。
\end{proof}

\begin{corollary}[有限 \(p\)-进精度下的 Gram--shift 有效分辨率]\label{cor:xi-gramshift-padic-effective-precision}
沿用命题 \ref{prop:xi-node-polynomial-good-reduction-discriminant} 的记号。
固定素数 \(p\)，并设 \((\widehat H_0,\widehat H_1)\in\Mat_\kappa(\ZZ_p)^2\) 为一组 \(p\)-进提升，使其在模 \(p^N\) 意义下与给定数据一致，且 \(\det(\widehat H_0)\ne 0\)。
令
\[
s:=v_p\!\bigl(\det(\widehat H_0)\bigr).
\]
则由同余数据 \((\widehat H_0,\widehat H_1)\bmod p^N\) 所确定的移位实现矩阵 \(\widehat A:=\widehat H_0^{-1}\widehat H_1\) 在 \(\Mat_\kappa(\ZZ_p/p^{N-s}\ZZ_p)\) 中是唯一的。
更具体地，若 \((\widehat H_0',\widehat H_1')\) 为另一组满足同一模 \(p^N\) 数据且 \(\det(\widehat H_0')\ne 0\) 的提升，并记 \(\widehat A':=\widehat H_0'^{-1}\widehat H_1'\)，则
\[
\widehat A\equiv \widehat A'\pmod{p^{N-s}}.
\]
此外，若在某一分裂扩张中存在 Prony 分解 \(s_m=\sum_{j=1}^{\kappa}\omega_j a_j^m\) 使
\(
\widehat H_0=(s_{i+j})_{0\le i,j\le \kappa-1}
\)，
则由推论 \ref{cor:xi-hankel-first-block-determinant-discriminant-weight} 得
\[
s=\sum_{j=1}^{\kappa}v_p(\omega_j)\;+\;v_p\!\left(\prod_{1\le i<j\le\kappa}(a_j-a_i)^2\right).
\]
\end{corollary}
\begin{proof}
由伴随矩阵公式
\(
\widehat H_0^{-1}=\operatorname{adj}(\widehat H_0)/\det(\widehat H_0)
\)，有
\[
\widehat A=\frac{\operatorname{adj}(\widehat H_0)\widehat H_1}{\det(\widehat H_0)}.
\]
注意 \(\operatorname{adj}(\widehat H_0)\widehat H_1\) 的每个条目均为 \(\widehat H_0,\widehat H_1\) 条目的整系数多项式，故同余条件
\(
(\widehat H_0,\widehat H_1)\equiv(\widehat H_0',\widehat H_1')\pmod{p^N}
\)
推出
\(
\operatorname{adj}(\widehat H_0)\widehat H_1\equiv \operatorname{adj}(\widehat H_0')\widehat H_1'\pmod{p^N}
\)。
再由 \(v_p(\det(\widehat H_0))=s\)（以及同余下 \(s=v_p(\det(\widehat H_0'))\)）可知分母提取出 \(p^s\) 后余因子为单位，故相除只损失 \(s\) 位精度，从而
\(
\widehat A\equiv \widehat A'\pmod{p^{N-s}}
\)。
末尾估值分解为推论 \ref{cor:xi-hankel-first-block-determinant-discriminant-weight} 的直接取赋值。证毕。
\end{proof}

\paragraph{\texorpdfstring{$p$}{p}-进有效分辨率与坏素数机制的进一步收紧}
在推论 \ref{cor:xi-gramshift-padic-effective-precision} 的“矩阵条目只损失 \(s=v_p(\det\widehat H_0)\) 位精度”之后，还可把节点层的可恢复精度与坏素数集合进一步压缩为两个完全内在的不变量：其一是特征多项式判别式的半指数损耗；其二是权重退化与节点碰撞所组成的二因子坏素数机制。

\begin{theorem}[\texorpdfstring{$p$}{p}-进有效分辨率的判别式半指数损耗律]\label{thm:xi-gramshift-padic-discriminant-halfloss}
\leanverified{paper\_xi\_gramshift\_padic\_discriminant\_halfloss}
沿用推论 \ref{cor:xi-gramshift-padic-effective-precision} 的记号。
设
\[
\widehat A:=\widehat H_0^{-1}\widehat H_1,\qquad
M:=N-s,\qquad s:=v_p(\det\widehat H_0),
\]
并令 \(\chi_{\widehat A}(T)=\det(TI_\kappa-\widehat A)\in\ZZ_p[T]\)。
取其某个分裂域 \(K/\QQ_p\)，并把 \(v_p\) 的延拓仍记作 \(v\)。
若 \(\chi_{\widehat A}\) 在 \(K\) 中分裂为互异根
\[
\chi_{\widehat A}(T)=\prod_{j=1}^{\kappa}(T-z_j),
\qquad
z_i\neq z_j\ (i\neq j),
\]
并记
\[
D:=v\!\bigl(\Disc(\chi_{\widehat A})\bigr),
\]
则从模 \(p^N\) 数据出发，节点多重集的最佳统一可恢复位阶满足
\begin{equation}\label{eq:xi-gramshift-padic-discriminant-halfloss-bound}
\boxed{\;
\mathfrak R_p(\chi_{\widehat A})
\ \le\
M-\Bigl\lceil\frac{D}{2}\Bigr\rceil,
\;}
\end{equation}
其中 \(\mathfrak R_p(\chi_{\widehat A})\) 表示“由该模 \(p^N\) 数据所能统一保证的最大 \(p\)-进位阶”。
等价地，不存在比
\[
v(z_j-z_j')\ \ge\ M-\Bigl\lceil\frac{D}{2}\Bigr\rceil
\qquad\text{（在适当重标号后）}
\]
更强的统一节点比较结论，其中 \(\{z_j'\}_{j=1}^{\kappa}\) 为任一与同一模 \(p^N\) 数据相容的另一组简单节点。
\par
进一步地，当 \(M\ge D+1\) 时，上述上界可由多变量 Hensel--Newton 反演达到；而在最坏情形下，损耗项 \(\lceil D/2\rceil\) 不能被任何统一估计改进。
\end{theorem}

\begin{proof}
由推论 \ref{cor:xi-gramshift-padic-effective-precision}，矩阵 \(\widehat A\) 在 \(\Mat_\kappa(\ZZ_p/p^{M}\ZZ_p)\) 中唯一确定，
故其特征多项式的系数向量也在模 \(p^{M}\) 意义下唯一确定。
\par
记
\[
\mathfrak e:(z_1,\dots,z_\kappa)\longmapsto
\bigl(e_1(z),\dots,e_\kappa(z)\bigr)
\]
为根向量到初等对称函数（等价地到首一特征多项式系数）的 Vieta 映射。
其雅可比行列式是经典 Vandermonde 行列式：
\[
\det D\mathfrak e(z)=\pm\prod_{1\le i<j\le\kappa}(z_j-z_i).
\]
因此
\[
v\!\bigl(\det D\mathfrak e(z)\bigr)
=
v\!\Bigl(\prod_{i<j}(z_j-z_i)\Bigr)
=
\frac12\,v\!\bigl(\Disc(\chi_{\widehat A})\bigr)
=
\frac{D}{2}.
\]
这说明在根层到系数层的局部线性化中，逆雅可比至少损失 \(D/2\) 个 \(p\)-进位阶；
故由系数精度 \(M\) 反演到根精度时，任何统一的局部恢复界都不可能优于
\(
M-\lceil D/2\rceil
\)，即得 \eqref{eq:xi-gramshift-padic-discriminant-halfloss-bound} 的上界性。
\par
若 \(M\ge D+1\)，则 \(\mathfrak e\) 在该简单根点满足多变量 Hensel--Newton 可逆条件，
并可按定理 \ref{thm:xi-nonarchimedean-prony-hensel-newton-inversion} 的同一证明机制得到局部唯一反演；
于是该上界在 Hensel 有效区间内可被达到。
\par
最后说明其最坏情形尖锐性。
在离散赋值域上对 \(D\mathfrak e(z)\) 取 Smith 正规形；
当全部坏赋值集中在单一不变量因子上（例如仅一对节点在赋值 \(m\) 上发生最近碰撞而其余节点彼此为单位分离，此时 \(D=2m\)）时，
\(K\) 的一致化元记为 \(\varpi\)，则
\(\mathfrak e\) 在局部解析坐标中可化为
\[
(x_1,\dots,x_\kappa)\longmapsto (\varpi^{m}x_1,x_2,\dots,x_\kappa)
\]
的同型模型。
于是系数层的 \(p\)-进精度 \(M\) 仅能把第一坐标确定到位阶 \(M-m\)，从而根层在位阶 \(M-m-1\) 处仍可发生可检测分歧。
这正给出最坏情形下的不可改进性。证毕。
\end{proof}

\begin{theorem}[有限指数分解下坏素数的二因子分解]\label{thm:xi-gramshift-bad-primes-weight-discriminant}
\leanverified{paper\_xi\_gramshift\_bad\_primes\_weight\_discriminant}
沿用定理 \ref{thm:xi-node-spectrum-minimal-window-2kappa} 与推论 \ref{cor:xi-node-polynomial-rational-reconstruction} 的记号。
设
\[
s_n=\sum_{j=1}^{\kappa}w_j z_j^n,
\qquad
H_0=(s_{i+j})_{0\le i,j\le \kappa-1},
\qquad
A:=H_0^{-1}H_1,
\qquad
\chi_A(T):=\det(TI_\kappa-A)=\prod_{j=1}^{\kappa}(T-z_j),
\]
其中 \(z_j\) 两两不同且 \(w_j\neq 0\)。
则有严格因式分解
\begin{equation}\label{eq:xi-gramshift-hankel-det-weight-disc-factorization}
\boxed{\;
\det(H_0)
=
\Bigl(\prod_{j=1}^{\kappa}w_j\Bigr)\cdot \Disc(\chi_A).
\;}
\end{equation}
若回到未归一化矩列 \(u_n=4\pi s_n\)，则对应 Hankel 指纹块 \(\mathsf H_{\kappa-1}=(u_{i+j})_{0\le i,j\le \kappa-1}\) 满足
\[
\det(\mathsf H_{\kappa-1})
=(4\pi)^\kappa
\Bigl(\prod_{j=1}^{\kappa}w_j\Bigr)\cdot \Disc(\chi_A).
\]
\par
进一步，固定素数 \(p\)，取包含全部 \(z_j,w_j\) 的有限扩张 \(K/\QQ_p\)，记其赋值环为 \(\mathcal O_K\)、极大理想为 \(\mathfrak p\)。
并取一组 \(p\)-整模型，使 \(H_0\)、\(\chi_A\) 与各权重 \(w_j\) 在 \(\mathcal O_K\) 中有意义。
则以下命题等价：
\begin{enumerate}
  \item \(H_0\bmod \mathfrak p\) 可逆；
  \item \(\prod_{j=1}^{\kappa}w_j\not\equiv 0\pmod{\mathfrak p}\) 且 \(\Disc(\chi_A)\not\equiv 0\pmod{\mathfrak p}\)；
  \item 模 \(\mathfrak p\) 意义下既无权重退化（某个 \(w_j\equiv 0\)），亦无节点碰撞（\(\chi_A\bmod \mathfrak p\) 在代数闭包中无重根）。
\end{enumerate}
因此，在有限指数分解模型中，坏素数集合严格分解为两类可解释机制的并：
\[
\{\text{坏素数}\}
=
\{\text{权重退化素数}\}\cup\{\text{节点碰撞素数}\}.
\]
结合定理 \ref{thm:xi-hankel-discriminant-unavoidable-bad-primes}，这说明该模型中的内在坏素数集合恰由积分模型下
\(
\bigl(\prod_j w_j\bigr)\Disc(\chi_A)
\)
的素因子支撑决定。
\end{theorem}

\begin{proof}
由 Hankel 的 Vandermonde 分解，
\[
H_0=V\,\diag(w_1,\dots,w_\kappa)\,V^{\mathsf T},
\qquad
V_{i,j}:=z_j^{\,i}\quad(0\le i\le \kappa-1),
\]
从而
\[
\det(H_0)=\Bigl(\prod_{j=1}^{\kappa}w_j\Bigr)\cdot (\det V)^2.
\]
而 \(A\) 的特征多项式为首一多项式 \(\chi_A(T)=\prod_j(T-z_j)\)，故
\[
\Disc(\chi_A)=\prod_{1\le i<j\le \kappa}(z_j-z_i)^2=(\det V)^2,
\]
即得 \eqref{eq:xi-gramshift-hankel-det-weight-disc-factorization}。
\par
在 \(p\)-整模型下，\(H_0\bmod \mathfrak p\) 可逆当且仅当 \(\det(H_0)\not\equiv 0\pmod{\mathfrak p}\)。
由上述分解，这等价于两因子同时为单位，即
\[
\prod_j w_j\not\equiv 0\pmod{\mathfrak p}
\qquad\text{且}\qquad
\Disc(\chi_A)\not\equiv 0\pmod{\mathfrak p}.
\]
第一条正是模 \(\mathfrak p\) 下无权重退化；
第二条则等价于 \(\chi_A\bmod \mathfrak p\) 为可分多项式，也即节点在代数闭包上无重合。
于是 (1)--(3) 等价。
\par
最后一段由定理 \ref{thm:xi-hankel-discriminant-unavoidable-bad-primes} 直接推出：坏素数集合既可由 Hankel 行列式含量内在描述，也可在此有限指数分解模型中解释为
\(
\bigl(\prod_j w_j\bigr)\Disc(\chi_A)
\)
的素因子支撑。证毕。
\end{proof}

\endinput


\paragraph{盘内极点参数与指数节点参数的闭合互译}
在有限点状缺陷—Toeplitz/扫描 Hankel 双接口的稳定区间内，同一缺陷原子既可由 Toeplitz 侧的盘内极点 \(a_j\in\mathbb{D}\) 表示，亦可由扫描侧的指数节点 \(z_j(\Delta)\in\mathbb{D}\) 表示（参见 \eqref{eq:xi-comoving-scan-fourier} 的指数谱闭式与 \eqref{eq:xi-scan-node-from-gamma-delta} 的节点定义）。
下面给出两种参数之间的无分支指数共轭恒等式，并据此把 Gram--shift 节点谱与 Toeplitz 盘几何严格对齐。

\begin{theorem}[Toeplitz 盘内极点与扫描指数节点之间的反全纯指数共轭]\label{thm:xi-toeplitz-diskpole-scan-node-antiholomorphic-conjugacy}
取 \(\Delta>0\)。
对任意盘点 \(a\in\mathbb{D}\)，令
\[
t(a):=\mathcal{C}^{-1}(a)=\mathrm{i}\frac{1+a}{1-a}\in\mathbb{H},
\qquad
t(a)=\gamma(a)+\mathrm{i}\delta(a)\quad(\delta(a)>0),
\]
并定义映射 \(\Phi_{\Delta}:\mathbb{D}\to\mathbb{D}\)
\begin{equation}\label{eq:xi-phi-diskpole-to-scan-node}
\boxed{\;
\Phi_{\Delta}(a)
:=e^{-\Delta}\exp\!\Bigl(-\Delta\,\overline{\frac{1+a}{1-a}}\Bigr).\;}
\end{equation}
则对任意缺陷原子 \((\gamma_j,\delta_j)\) 的盘内极点
\(
a_j:=\mathcal{C}(\gamma_j+\mathrm{i}\delta_j)\in\mathbb{D}
\)
与同一步长下的扫描指数节点
\begin{equation}\label{eq:xi-scan-node-from-gamma-delta}
z_j(\Delta):=\exp\!\bigl(-(1+\delta_j)\Delta\bigr)\exp\!\bigl(-\mathrm{i}\gamma_j\Delta\bigr)\in\mathbb{D},
\end{equation}
恒有无分支歧义的恒等式
\[
\boxed{\;z_j(\Delta)=\Phi_{\Delta}(a_j).\;}
\]
\end{theorem}
\begin{proof}
由 \(a_j=\mathcal{C}(\gamma_j+\mathrm{i}\delta_j)\) 与 \(t(a)=\mathcal{C}^{-1}(a)\) 的定义，得
\(
t(a_j)=\gamma_j+\mathrm{i}\delta_j
\)。
记 \(t_j:=t(a_j)\)。
由
\(
t_j=\mathrm{i}\frac{1+a_j}{1-a_j}
\)
取共轭并乘以 \(\mathrm{i}\) 得
\[
\mathrm{i}\,\overline{t_j}
\;=\;
\mathrm{i}\,\overline{\mathrm{i}\frac{1+a_j}{1-a_j}}
\;=\;
\overline{\frac{1+a_j}{1-a_j}}.
\]
另一方面，\(\mathrm{i}\overline{t_j}=\mathrm{i}(\gamma_j-\mathrm{i}\delta_j)=\delta_j+\mathrm{i}\gamma_j\)，故
\(
\delta_j+\mathrm{i}\gamma_j=\overline{\frac{1+a_j}{1-a_j}}
\)。
将其代回 \eqref{eq:xi-scan-node-from-gamma-delta} 可得
\[
z_j(\Delta)
=e^{-\Delta}\exp\!\bigl(-\Delta(\delta_j+\mathrm{i}\gamma_j)\bigr)
=e^{-\Delta}\exp\!\Bigl(-\Delta\,\overline{\frac{1+a_j}{1-a_j}}\Bigr)
=\Phi_\Delta(a_j).
\]
证毕。
\end{proof}

\begin{proposition}[两点比值的端点归一化 Cayley--指数公式]\label{prop:xi-phi-diskpole-ratio-identity}
\leanverified{paper\_xi\_phi\_diskpole\_ratio\_identity}
沿定理 \ref{thm:xi-toeplitz-diskpole-scan-node-antiholomorphic-conjugacy} 的记号。
对任意 \(a,b\in\mathbb{D}\) 有严格恒等式
\begin{equation}\label{eq:xi-phi-ratio-identity}
\boxed{\;
\frac{\Phi_{\Delta}(a)}{\Phi_{\Delta}(b)}
=\exp\!\Bigl(-2\Delta\,\overline{\frac{a-b}{(1-a)(1-b)}}\Bigr).\;}
\end{equation}
右端为整函数复合，故不涉及复对数的分支选择。
\end{proposition}
\begin{proof}
直接计算得
\[
\frac{1+a}{1-a}-\frac{1+b}{1-b}
=\frac{(1+a)(1-b)-(1+b)(1-a)}{(1-a)(1-b)}
=\frac{2(a-b)}{(1-a)(1-b)}.
\]
代入 \eqref{eq:xi-phi-diskpole-to-scan-node} 即得 \eqref{eq:xi-phi-ratio-identity}。证毕。
\end{proof}

\begin{corollary}[Gram--shift 节点谱作为 Toeplitz 盘内极点集的 \texorpdfstring{$\Phi_{\Delta}$}{Phi}-像]\label{cor:xi-gramshift-spectrum-as-phi-image-of-diskpoles}
在定理 \ref{thm:xi-gram-shift-positive-pencil-node-spectrum} 的稳定区间设定下，
假设点状缺陷由有限盘内极点多重集 \(\mathcal A=\{a_1,\dots,a_\kappa\}\subset\mathbb{D}\) 承载，
并且扫描节点 \(\{z_j(\Delta)\}_{j=1}^{\kappa}\) 与盘内极点 \(\{a_j\}\) 由同一组缺陷原子 \((\gamma_j,\delta_j)\) 诱导（即 \eqref{eq:xi-scan-node-from-gamma-delta} 与 \(a_j=\mathcal{C}(\gamma_j+\mathrm{i}\delta_j)\) 同时成立）。
则 Gram--shift 实现矩阵 \(A:=H_0^{-1}H_1\) 的谱满足
\[
\sigma(A)=\Phi_\Delta(\mathcal A):=\{\Phi_\Delta(a_1),\dots,\Phi_\Delta(a_\kappa)\},
\]
其中按代数重数计。
\end{corollary}
\leanverified{paper\_xi\_gramshift\_spectrum\_as\_phi\_image\_of\_diskpoles}
\begin{proof}
由定理 \ref{thm:xi-gram-shift-positive-pencil-node-spectrum} 得 \(\sigma(A)=\{z_j(\Delta)\}_{j=1}^{\kappa}\)。
再由定理 \ref{thm:xi-toeplitz-diskpole-scan-node-antiholomorphic-conjugacy}，对每个 \(j\) 有 \(z_j(\Delta)=\Phi_\Delta(a_j)\)。
逐点对应并按重数计即可。证毕。
\end{proof}

\begin{proposition}[Gram--shift 特征多项式的 Toeplitz--几何闭式表示]\label{prop:xi-gramshift-charpoly-as-phi-image-of-diskpoles}
在推论 \ref{cor:xi-gramshift-spectrum-as-phi-image-of-diskpoles} 的设定下，
Gram--shift 矩阵 \(A\) 的特征多项式满足
\[
\chi_A(\lambda)=\prod_{j=1}^{\kappa}\bigl(\lambda-\Phi_\Delta(a_j)\bigr),
\]
其中 \(\mathcal A=\{a_j\}\) 按重数计。
\end{proposition}
\begin{proof}
由推论 \ref{cor:xi-gramshift-spectrum-as-phi-image-of-diskpoles}，\(A\) 的全部特征值（按重数）为 \(\{\Phi_\Delta(a_j)\}\)，故特征多项式为所示乘积。证毕。
\end{proof}

\begin{proposition}[扫描谱指纹的盘几何反全纯表达]\label{prop:xi-scan-fingerprint-antiholomorphic-functional-of-diskpoles}
沿用定理 \ref{thm:xi-comoving-scan-hankel-rank-defect} 的记号，取步长 \(\Delta>0\) 并令 \(u_n:=\widehat H(n\Delta)\)（\(n\ge 0\)）。
在有限点状缺陷口径下，设盘内极点为 \(\mathcal A=\{a_1,\dots,a_\kappa\}\subset\mathbb{D}\)，并记
\[
t_j:=\mathcal{C}^{-1}(a_j)=\gamma_j+\mathrm{i}\delta_j\in\mathbb{H},
\qquad
\delta_j>0,
\]
且每个原子携带权 \(m_j\in\NN\)（与 \eqref{eq:xi-comoving-scan-fourier} 的重数记号一致）。
则对所有 \(n\ge 0\) 成立闭式恒等式
\[
u_n
=4\pi\sum_{j=1}^{\kappa}m_j\,\frac{\delta_j}{1+\delta_j}\,\bigl(\Phi_\Delta(a_j)\bigr)^{n}.
\]
\end{proposition}
\begin{proof}
由 \eqref{eq:xi-comoving-scan-fourier} 与 \(u_n=\widehat H(n\Delta)\) 得
\[
u_n
=4\pi\sum_{j=1}^{\kappa}m_j\,\frac{\delta_j}{1+\delta_j}\,e^{-(1+\delta_j)n\Delta}\,e^{-\mathrm{i}\gamma_j n\Delta}.
\]
而由定理 \ref{thm:xi-toeplitz-diskpole-scan-node-antiholomorphic-conjugacy}，
\(
\Phi_\Delta(a_j)=e^{-(1+\delta_j)\Delta}e^{-\mathrm{i}\gamma_j\Delta}
\)，
故 \((\Phi_\Delta(a_j))^{n}=e^{-(1+\delta_j)n\Delta}e^{-\mathrm{i}\gamma_j n\Delta}\)。
代回即得所示恒等式。证毕。
\end{proof}

\begin{corollary}[节点谱的噪声稳定性：矩阵差到谱差再到对数参数差]\label{cor:xi-node-spectrum-noise-stability}
沿用定理 \ref{thm:xi-gram-shift-positive-pencil-node-spectrum} 的记号。
给定噪声观测 \(\tilde u_n=u_n+\delta_n\) 并据此构造 \(\widetilde H_0,\widetilde H_1\) 与 \(\widetilde A:=\widetilde H_0^{-1}\widetilde H_1\)。
若满足小扰动门槛
\[
\|H_0^{-1}\|\,\|\widetilde H_0-H_0\|\le \frac12,
\]
则由命题 \ref{prop:xi-hankel-inversion-noise-stability-neumann-bauerfike} 得到显式扰动界
\[
\|\widetilde A-A\|
\le 4\,\|H_0^{-1}\|\Bigl(\|\widetilde H_1-H_1\|+\|A\|\,\|\widetilde H_0-H_0\|\Bigr),
\]
并且当 \(A\) 可对角化 \(A=PDP^{-1}\) 时，任意 \(\widetilde A\) 的特征值 \(\widetilde\lambda\) 满足 Bauer--Fike 型谱稳定界
\[
\mathrm{dist}\bigl(\widetilde\lambda,\sigma(A)\bigr)\le \kappa(P)\,\|\widetilde A-A\|,
\qquad \kappa(P):=\|P\|\,\|P^{-1}\|.
\]
若另有谱下界 \(\inf_{\lambda\in\sigma(A)}|\lambda|\ge r>0\) 且 \(|\widetilde\lambda-\lambda|\le r/2\)，则对数参数满足
\[
|\log\widetilde\lambda-\log\lambda|\le \frac{2}{r}\,|\widetilde\lambda-\lambda|.
\]
\end{corollary}
\begin{proof}
前两条由命题 \ref{prop:xi-hankel-inversion-noise-stability-neumann-bauerfike} 将 \(H_0,H_1\) 专门化为本段的移位 Hankel 窗口即得。
最后一条由解析函数估计得到：在闭环域 \(\{z:|z|\ge r/2\}\) 上有 \(|(\log z)'|=1/|z|\le 2/r\)，故 \(\log\) 在该域上为 Lipschitz。证毕。
\end{proof}

\begin{theorem}[Gram--shift 节点谱与 Toeplitz 铅笔极点谱的 Cayley--log 互译]\label{thm:xi-gramshift-toeplitz-cayley-log-bridge}
在定理 \ref{thm:xi-gram-shift-positive-pencil-node-spectrum} 的有限缺陷设定下，假设扫描谱指纹满足指数分解 \eqref{eq:xi-scan-fingerprint-exponential-sum}，并进一步满足节点可参数化为
\[
z_j=e^{-(1+\delta_j)}e^{-\mathrm{i}\gamma_j}\qquad(1\le j\le\kappa),
\]
其中 \(\delta_j>0\)、\(\gamma_j\in\RR\)。
令 Cayley 紧化 \(\mathcal{C}:\mathbb{H}\to\mathbb{D}\) 如 \S\ref{subsubsec:xi-pontryagin-toeplitz-blaschke-tomography-tate}，并定义真实盘内极点
\[
a_j:=\mathcal{C}(\gamma_j+\mathrm{i}\delta_j)\in\mathbb{D}.
\]
在有限型且负惯性指标在高阶截断上稳定的情形，Toeplitz 铅笔 \((\mathbf{T}_N,\mathbf{T}_N^{(1)})\) 的盘内广义谱
\(
\Lambda_N:=\mathrm{Spec}(\mathbf{T}_N,\mathbf{T}_N^{(1)})\cap\mathbb{D}
\)
在 \(N\to\infty\) 时收敛到该极点多重集合（结论 \ref{con:xi-finite-falsifier-offline-fiber-readout}）。
取主值辐角 \(\mathrm{Arg}(z)\in(-\pi,\pi]\)，并对满足 \(-\log|z|>1\) 的 \(z\in\mathbb{D}\setminus\{0\}\) 定义
\[
\Phi(z):=\mathcal{C}\!\Bigl(-\mathrm{Arg}(z)+\mathrm{i}\bigl(-\log|z|-1\bigr)\Bigr).
\]
则对每个 \(j\) 成立严格恒等式
\[
\boxed{\ a_j=\Phi(z_j)\ }.
\]
因此 Gram--shift 节点谱与 Toeplitz 铅笔极点谱在盘内极点层面由 \(\Phi\) 显式对接。
\end{theorem}
\begin{proof}
由节点参数化直接得到
\(
\gamma_j=-\mathrm{Arg}(z_j)
\)
与
\(
\delta_j=-\log|z_j|-1
\)（此处 \(\delta_j>0\) 等价于 \(-\log|z_j|>1\)）。
代入 \(a_j=\mathcal{C}(\gamma_j+\mathrm{i}\delta_j)\) 即得 \(a_j=\Phi(z_j)\)。证毕。
\end{proof}
\leanverified{paper\_xi\_gramshift\_toeplitz\_cayley\_log\_bridge}

\begin{theorem}[\texorpdfstring{$2\kappa$}{2k} 样本全息：由有限扫描样本重建负谱体积势]\label{thm:xi-2k-sample-holography-reconstructs-detP}
在定理 \ref{thm:xi-gramshift-toeplitz-cayley-log-bridge} 的设定下，设 \(\kappa\ge 1\) 并假设进入稳定区间使得 \(H_0=\mathsf H_{\kappa-1}\) 可逆（定理 \ref{thm:xi-gram-shift-positive-pencil-node-spectrum}）。
则有限样本族
\[
\{u_n\}_{0\le n\le 2\kappa-1}
\]
唯一确定如下对象：
\begin{enumerate}
\normalfont
  \item 缺陷节点谱 \(\{z_j\}_{j=1}^{\kappa}\subset\mathbb{D}\)（按代数重数计）；
  \item 盘内极点多重集合 \(\{a_j\}_{j=1}^{\kappa}\subset\mathbb{D}\)（按代数重数计）；
  \item Pick--Poisson Gram 矩阵 \(\mathsf P=\mathsf P(\{a_j\})\) 的行列式 \(\det\mathsf P\) 以及负谱体积势
  \[
  S_{\mathrm{gen}}:=-\log\det\mathsf P.
  \]
\end{enumerate}
此外，在 Toeplitz 主子矩阵族进入稳定显影区间且负惯性最终稳定为 \(\kappa\) 的情形，若记其负特征值（按重数）为 \(\{\lambda_j^{-}(\mathbf{T}_N)\}_{j=1}^{\kappa}\)，则满足乘积极限
\[
\lim_{N\to\infty}\ \prod_{j=1}^{\kappa}\bigl(-\lambda_j^{-}(\mathbf{T}_N)\bigr)=\det\mathsf P
\]
（结论 \ref{con:xi-terminal-negative-spectrum-det-invariant}）。
因此 Toeplitz 负谱“体积极限”亦由 \(\{u_n\}_{0\le n\le 2\kappa-1}\) 全息决定。
\end{theorem}
\begin{proof}
由定理 \ref{thm:xi-gram-shift-positive-pencil-node-spectrum}，样本 \(\{u_n\}_{0\le n\le 2\kappa-1}\) 唯一确定移位 Hankel 块 \(H_0,H_1\)，从而确定矩阵 \(A=H_0^{-1}H_1\) 并由 \(\sigma(A)=\{z_j\}\) 得到节点谱。
再由定理 \ref{thm:xi-gramshift-toeplitz-cayley-log-bridge} 得到 \(\{a_j\}\) 的代数确定性。
最后，\(\det\mathsf P\) 由定理 \ref{thm:xi-pick-poisson-vandermonde-resultant-identity}（或等价的命题 \ref{prop:xi-pick-poisson-det-factorization}）作为 \(\{a_j\}\) 的函数完全确定，从而 \(S_{\mathrm{gen}}=-\log\det\mathsf P\) 亦被唯一确定。证毕。
\end{proof}
\leanverified{paper\_xi\_2k\_sample\_holography\_reconstructs\_detp}

\begin{conjecture}[稳定显影区间内 Toeplitz 铅笔与 Gram--shift 铅笔的有限阶谱刚性]\label{conj:xi-toeplitz-gramshift-finite-order-spectral-rigidity}
在定理 \ref{thm:xi-gramshift-toeplitz-cayley-log-bridge} 的有限型与稳定显影区间设定下，
存在 \(N_0\) 使对所有 \(N\ge N_0\) 都有多重集恒等式
\[
\boxed{\;
\Lambda_N=\Phi\!\bigl(\sigma(A)\bigr)\;}
\qquad
\Bigl(
\Lambda_N:=\mathrm{Spec}(\mathbf{T}_N,\mathbf{T}_N^{(1)})\cap\mathbb{D},\ 
A:=H_0^{-1}H_1
\Bigr),
\]
其中 \(\Phi\) 由定理 \ref{thm:xi-gramshift-toeplitz-cayley-log-bridge} 中的定义给出。
换言之，Toeplitz 铅笔的盘内广义特征谱在充分大阶数后应当发生\emph{终端冻结}，并与扫描指纹的 Gram--shift 节点谱在同一坐标图下完全一致（而非仅在 \(N\to\infty\) 意义下收敛）。
\end{conjecture}

\begin{remark}[支持性论证（非证明）]
该猜想与正文中两条独立机制在稳定显影区间内的对齐相容：
\begin{enumerate}
  \item \textnormal{（Toeplitz 侧的有限否证器收敛）}结论 \ref{con:xi-finite-falsifier-offline-fiber-readout} 给出 \(\Lambda_N\) 在 \(N\to\infty\) 时收敛到真实盘内极点多重集的有限型口径；
  若进一步假设稳定后不存在“赝谱候选”（即盘内候选点均来自真实极点），则自然期望收敛在有限阶后终止为精确一致。
  \item \textnormal{（Gram--shift 侧的有限窗精确实现）}定理 \ref{thm:xi-gram-shift-positive-pencil-node-spectrum} 在有限指数分解下给出
  \(\sigma(A)=\{z_j\}\) 的严格代数实现；
  同时定理 \ref{thm:xi-gramshift-toeplitz-cayley-log-bridge} 给出节点与极点的显式互译 \(\Phi\)，从而把两类谱读出放入同一盘内坐标系比较。
  \item \textnormal{（负方向稳定的刚性暗示）}在稳定显影区间内，Toeplitz 负方向由 Hankel 指纹块精确刻画且在可允许扰动半径内保持不变（定理 \ref{thm:xi-toeplitz-negative-inertia-spectral-gap-stability}）。
  该“稳定负子空间已承载全部点状缺陷数据”的结构，为有限阶谱冻结提供了额外刚性来源。
\end{enumerate}
\end{remark}


\begin{proposition}[Hankel 满秩的一素数审计]\label{prop:xi-hankel-fullrank-one-prime-audit}
设 \(H_0\in \Mat_\kappa(\QQ)\)。
取公共分母 \(D\in\NN\) 使 \(\widehat H_0:=DH_0\in \Mat_\kappa(\ZZ)\)。
则：
\begin{enumerate}
  \item \(H_0\) 在 \(\QQ\) 上可逆当且仅当 \(\det(\widehat H_0)\neq 0\)（整数意义）。
  \item 若存在素数 \(p\nmid D\) 使 \(\det(\widehat H_0)\not\equiv 0\ (\mathrm{mod}\ p)\)，则 \(\det(\widehat H_0)\neq 0\)，从而 \(H_0\) 在 \(\QQ\) 上可逆。
  \item 若 \(\det(\widehat H_0)\neq 0\)，则满足 \(\det(\widehat H_0)\equiv 0\ (\mathrm{mod}\ p)\) 的素数集合有限，且必包含于 \(\det(\widehat H_0)\) 的素因子集合。
\end{enumerate}
\leanverified{paper\_xi\_hankel\_fullrank\_one\_prime\_audit}
\end{proposition}
\begin{proof}
第一条由 \(\det(\widehat H_0)=D^\kappa\det(H_0)\) 立得。
第二条：若 \(\det(\widehat H_0)=0\) 作为整数成立，则其模任意素数皆为零，与假设矛盾。
第三条：若 \(p\mid \det(\widehat H_0)\)，则必有 \(\det(\widehat H_0)\equiv 0\ (\mathrm{mod}\ p)\)；反之模零亦推出整除。
由于非零整数仅有有限个素因子，故坏素数集合有限。证毕。
\end{proof}

\leanverified{paper\_xi\_hankel\_linear\_realization\_finite\_prime\_support\_mod\_audit}
\begin{proposition}[Hankel 线性实现的有限素数支撑与模审计验证器]\label{prop:xi-hankel-linear-realization-finite-prime-support-mod-audit}
设 \(\widehat H_0,\widehat H_1\in\Mat_{\kappa}(\ZZ)\) 且 \(\det(\widehat H_0)\neq 0\)，并定义有理实现矩阵
\[
A:=\widehat H_0^{-1}\widehat H_1\in\Mat_{\kappa}(\QQ).
\]
则：
\begin{enumerate}
  \item \textnormal{（有限素数支撑）}对任意素数 \(p\nmid\det(\widehat H_0)\) 与任意 \(k\ge 1\)，线性同余
  \[
  \widehat H_0X\equiv \widehat H_1\pmod{p^k}
  \]
  在 \(\Mat_{\kappa}(\ZZ/p^k\ZZ)\) 中存在唯一解，并且该解满足 \(X\equiv A\pmod{p^k}\)。
  因而所有实现歧义只能发生在有限集合 \(S:=\{p:\ p\mid \det(\widehat H_0)\}\) 上。
  \item \textnormal{（压缩模审计）}给定候选 \(\widehat A\in\Mat_{\kappa}(\QQ)\)，取 \(N^\ast\in\NN\) 使 \(N^\ast\widehat A\in\Mat_{\kappa}(\ZZ)\)，并定义
  \[
  F:=N^\ast(\widehat H_0\widehat A-\widehat H_1)\in\Mat_{\kappa}(\ZZ),
  \qquad
  |F|_\infty:=\max_{i,j}|F_{ij}|.
  \]
  若取有限素数集 \(\mathcal P\) 使 \(\prod_{p\in\mathcal P}p>2|F|_\infty\)，并且满足
  \[
  F\equiv 0\pmod p\qquad(\forall p\in\mathcal P),
  \]
  则必有 \(F=0\)，从而 \(\widehat H_0\widehat A=\widehat H_1\) 并推出 \(\widehat A=A\)。
\end{enumerate}
\end{proposition}

\begin{proof}
第一条：若 \(p\nmid\det(\widehat H_0)\)，则 \(\widehat H_0\) 在 \(\ZZ/p^k\ZZ\) 上可逆，故同余方程存在唯一解
\(
X\equiv \widehat H_0^{-1}\widehat H_1\equiv A\ (\mathrm{mod}\ p^k)
\)。
\par
第二条：若对所有 \(p\in\mathcal P\) 有 \(F\equiv 0\ (\mathrm{mod}\ p)\)，则 \(F\) 的每个条目都被 \(\prod_{p\in\mathcal P}p\) 整除。
由 \(\prod_{p\in\mathcal P}p>2|F|_\infty\) 得每个条目只能为 \(0\)，故 \(F=0\)。
于是 \(\widehat H_0\widehat A=\widehat H_1\)，再由 \(\widehat H_0\) 在 \(\QQ\) 上可逆得 \(\widehat A=A\)。证毕。
\end{proof}

\begin{corollary}[缺陷熵迫使负块谱质量非退化]\label{cor:xi-defect-entropy-forces-negative-block-mass}
在推论 \ref{cor:xi-toeplitz-scan-hankel-congruence-decomposition} 的稳定区间内，负块为
\(
G_{\kappa-1}:=\mathsf H_{\kappa-1}^\ast\mathsf H_{\kappa-1}\succeq 0
\)。
则其算子范数满足严格下界
\[
\|G_{\kappa-1}\|_{\mathrm{op}}
\ \ge\
\frac{|u_0|^2}{\kappa}.
\]
并且由命题 \ref{prop:xi-defect-entropy-hyperbolic-area-law-4pi} 的 \(4\pi\) 规一化，
\(
u_0=\widehat H(0)=\int_{\RR}H(x)\,\dd x=4\pi S_{\mathrm{def}}
\)，从而
\[
\|G_{\kappa-1}\|_{\mathrm{op}}\ \ge\ \frac{(4\pi S_{\mathrm{def}})^2}{\kappa}.
\]
\end{corollary}
\begin{proof}
由 \(G_{\kappa-1}=\mathsf H_{\kappa-1}^\ast\mathsf H_{\kappa-1}\) 得
\(
\|G_{\kappa-1}\|_{\mathrm{op}}=\sigma_{\max}(\mathsf H_{\kappa-1})^2
\)。
又由 Frobenius 范数与奇异值关系，
\[
\sigma_{\max}(\mathsf H_{\kappa-1})^2
\ge \frac{1}{\kappa}\sum_{i=1}^{\kappa}\sigma_i(\mathsf H_{\kappa-1})^2
=\frac{\|\mathsf H_{\kappa-1}\|_F^2}{\kappa}
\ge \frac{|(\mathsf H_{\kappa-1})_{1,1}|^2}{\kappa}
=\frac{|u_0|^2}{\kappa},
\]
得到第一条不等式。
最后将 \(u_0=4\pi S_{\mathrm{def}}\) 代入即得第二条不等式。证毕。
\end{proof}

% Entropy gap suppression and Hankel collapse bounds.

\begin{theorem}[熵缺口对非零谱指纹的指数抑制]\label{thm:xi-entropy-gap-exponential-suppression-nonzero-fingerprint}
沿用 \eqref{eq:xi-comoving-scan-fourier} 的有限缺陷口径，并取采样步长 \(\Delta=1\)，令
\(
u_n:=\widehat H(n)
\)
为离散谱指纹序列。
记整数缺陷指数 \(\kappa_{\mathrm{ind}}:=\sum_{j=1}^{\kappa}m_j\) 与缺陷熵 \(S_{\mathrm{def}}\)（\eqref{eq:xi-defect-entropy-closed-form}），并定义熵缺口
\[
\kappa_{\mathrm{ind}}-S_{\mathrm{def}}
=\sum_{j=1}^{\kappa}\frac{m_j}{1+\delta_j}
\qquad\text{（见 \eqref{eq:xi-defect-entropy-index-decomposition}）}.
\]
则对每个整数 \(n\ge 1\) 成立显式上界
\[
|u_n|
\ \le\
4\pi\,e^{-n}\,(\kappa_{\mathrm{ind}}-S_{\mathrm{def}}).
\]
\end{theorem}
\begin{proof}
由 \eqref{eq:xi-comoving-scan-fourier}（取 \(\Delta=1\)）可写
\[
u_n
=4\pi\sum_{j=1}^{\kappa}m_j\,\frac{\delta_j}{1+\delta_j}\,e^{-(1+\delta_j)n}\,e^{-\mathrm{i}\gamma_j n}.
\]
取模并分离指数因子得
\[
|u_n|
\le
4\pi\,e^{-n}\sum_{j=1}^{\kappa}m_j\,\frac{\delta_j}{1+\delta_j}\,(e^{-\delta_j})^{n}.
\]
对任意 \(\delta>0\) 有 \(e^\delta\ge 1+\delta\)，从而 \(e^{-\delta}\le (1+\delta)^{-1}\)；
于是对 \(n\ge 1\),
\(
(e^{-\delta_j})^{n}\le (1+\delta_j)^{-n}\le (1+\delta_j)^{-1}
\)。
代回并用 \(\frac{\delta_j}{1+\delta_j}\le 1\) 得
\[
|u_n|
\le
4\pi\,e^{-n}\sum_{j=1}^{\kappa}m_j\,\frac{1}{1+\delta_j}
=
4\pi\,e^{-n}\,(\kappa_{\mathrm{ind}}-S_{\mathrm{def}}),
\]
其中最后一步使用 \eqref{eq:xi-defect-entropy-index-decomposition}。证毕。
\end{proof}
\leanverified{paper\_xi\_entropy\_gap\_exponential\_suppression\_nonzero\_fingerprint}

\begin{proposition}[非零谐波的熵乘积锐化界]\label{prop:xi-nonzero-harmonic-entropy-gap-product-sharp}
在定理 \ref{thm:xi-entropy-gap-exponential-suppression-nonzero-fingerprint} 的设定下，令
\[
w_j:=\frac{\delta_j}{1+\delta_j}\in(0,1),
\qquad
S_{\mathrm{def}}=\sum_{j=1}^{\kappa}m_j w_j,
\qquad
\kappa_{\mathrm{ind}}-S_{\mathrm{def}}=\sum_{j=1}^{\kappa}m_j(1-w_j).
\]
则对每个整数 \(n\ge 1\) 成立统一上界
\[
|u_n|
\ \le\
4\pi\,e^{-n}\,\frac{S_{\mathrm{def}}(\kappa_{\mathrm{ind}}-S_{\mathrm{def}})}{\kappa_{\mathrm{ind}}}.
\]
\end{proposition}
\begin{proof}
由 \eqref{eq:xi-comoving-scan-fourier}（取 \(\Delta=1\)）并引入 \(w_j=\delta_j/(1+\delta_j)\) 得
\[
u_n
=4\pi\,e^{-n}\sum_{j=1}^{\kappa}m_j\,w_j\,e^{-\delta_j n}\,e^{-\mathrm{i}\gamma_j n}.
\]
由 \(e^\delta\ge 1+\delta\) 推出 \(e^{-\delta}\le (1+\delta)^{-1}=1-w\)，从而对 \(n\ge 1\) 有
\(
e^{-\delta_j n}\le (1-w_j)^n\le (1-w_j)
\)。
因此
\[
|u_n|
\le
4\pi\,e^{-n}\sum_{j=1}^{\kappa}m_j\,w_j(1-w_j)
=
4\pi\,e^{-n}\Bigl(S_{\mathrm{def}}-\sum_{j=1}^{\kappa}m_j w_j^2\Bigr).
\]
将 \((w_j)\) 按重数展开为长度 \(\kappa_{\mathrm{ind}}\) 的序列并应用 Cauchy--Schwarz 不等式，
\[
\sum_{j=1}^{\kappa}m_j w_j^2
\ \ge\
\frac{\bigl(\sum_{j=1}^{\kappa}m_j w_j\bigr)^2}{\sum_{j=1}^{\kappa}m_j}
=\frac{S_{\mathrm{def}}^2}{\kappa_{\mathrm{ind}}}.
\]
代回即得结论。证毕。
\end{proof}
\leanverified{paper\_xi\_nonzero\_harmonic\_entropy\_gap\_product\_sharp}

\begin{corollary}[两点采样反演的熵缺口下界]\label{cor:xi-entropy-gap-lower-bound-from-two-samples}
在命题 \ref{prop:xi-nonzero-harmonic-entropy-gap-product-sharp} 的设定下，若 \(u_0=4\pi S_{\mathrm{def}}\)，则对每个整数 \(n\ge 1\) 有
\[
\kappa_{\mathrm{ind}}-S_{\mathrm{def}}
\ \ge\
\kappa_{\mathrm{ind}}\,e^{n}\,\frac{|u_n|}{|u_0|}.
\]
\end{corollary}
\begin{proof}
由命题 \ref{prop:xi-nonzero-harmonic-entropy-gap-product-sharp} 得
\[
|u_n|
\le
4\pi\,e^{-n}\,\frac{S_{\mathrm{def}}(\kappa_{\mathrm{ind}}-S_{\mathrm{def}})}{\kappa_{\mathrm{ind}}}.
\]
结合 \(u_0=4\pi S_{\mathrm{def}}\) 并整理即得。证毕。
\end{proof}
\leanverified{paper\_xi\_entropy\_gap\_lower\_bound\_from\_two\_samples}

\begin{proposition}[第一谐波的方差惩罚律]\label{prop:xi-first-harmonic-variance-penalty}
\leanverified{paper\_xi\_first\_harmonic\_variance\_penalty}
在命题 \ref{prop:xi-nonzero-harmonic-entropy-gap-product-sharp} 的设定下，将 \((w_j)\) 按重数展开为长度 \(\kappa_{\mathrm{ind}}\) 的序列 \((w_\ell)_{\ell=1}^{\kappa_{\mathrm{ind}}}\subset(0,1)\) 并记均值
\[
\bar w:=\frac{1}{\kappa_{\mathrm{ind}}}\sum_{\ell=1}^{\kappa_{\mathrm{ind}}}w_\ell=\frac{S_{\mathrm{def}}}{\kappa_{\mathrm{ind}}},
\qquad
\Var(w):=\frac{1}{\kappa_{\mathrm{ind}}}\sum_{\ell=1}^{\kappa_{\mathrm{ind}}}(w_\ell-\bar w)^2.
\]
则第一谐波满足显式上界
\[
\frac{e}{4\pi}\,|u_1|
\ \le\
\frac{S_{\mathrm{def}}(\kappa_{\mathrm{ind}}-S_{\mathrm{def}})}{\kappa_{\mathrm{ind}}}
\ -\
\kappa_{\mathrm{ind}}\,\Var(w).
\]
\end{proposition}
\begin{proof}
由 \eqref{eq:xi-comoving-scan-fourier}（取 \(\Delta=1\) 且 \(n=1\)）并沿用 \(e^{-\delta}\le (1+\delta)^{-1}=1-w\) 得
\[
|u_1|
\le
4\pi\,e^{-1}\sum_{\ell=1}^{\kappa_{\mathrm{ind}}}w_\ell(1-w_\ell)
=
4\pi\,e^{-1}\Bigl(\sum_{\ell=1}^{\kappa_{\mathrm{ind}}}w_\ell-\sum_{\ell=1}^{\kappa_{\mathrm{ind}}}w_\ell^2\Bigr).
\]
注意
\[
\sum_{\ell=1}^{\kappa_{\mathrm{ind}}}w_\ell^2
=
\kappa_{\mathrm{ind}}\Var(w)+\kappa_{\mathrm{ind}}\bar w^2
=
\kappa_{\mathrm{ind}}\Var(w)+\frac{S_{\mathrm{def}}^2}{\kappa_{\mathrm{ind}}},
\]
代回并整理即得结论。证毕。
\end{proof}

\begin{proposition}[熵缺口驱动的 Hankel 主子式塌缩上界]\label{prop:xi-entropy-gap-hankel-determinant-collapse-upper-bound}
\leanverified{paper\_xi\_entropy\_gap\_hankel\_determinant\_collapse\_upper\_bound}
令 \(\mathsf H_{\kappa-1}:=(u_{p+q})_{0\le p,q\le \kappa-1}\) 为 Toeplitz--扫描 Hankel 框架中的 Hankel 指纹块，并定义
\[
B:=\sup_{n\ge 1}|u_n|.
\]
则行列式满足组合学型上界
\[
|\det(\mathsf H_{\kappa-1})|
\ \le\
(\kappa-1)!\,|u_0|\,B^{\kappa-1}+\bigl(\kappa!-(\kappa-1)!\bigr)\,B^{\kappa}
\ \le\
\kappa!\bigl(|u_0|\,B^{\kappa-1}+B^{\kappa}\bigr).
\]
特别地，在定理 \ref{thm:xi-entropy-gap-exponential-suppression-nonzero-fingerprint} 的设定下，若 \(S_{\mathrm{def}}\) 固定而 \(\kappa_{\mathrm{ind}}-S_{\mathrm{def}}\to 0\)，则存在仅依赖 \(\kappa\) 的常数使
\[
|\det(\mathsf H_{\kappa-1})|\ \ll_{\kappa}\ (\kappa_{\mathrm{ind}}-S_{\mathrm{def}})^{\kappa-1}.
\]
\end{proposition}
\begin{proof}
以排列展开式写
\[
\det(\mathsf H_{\kappa-1})
=\sum_{\pi\in S_{\kappa}}\mathrm{sgn}(\pi)\prod_{p=0}^{\kappa-1}u_{p+\pi(p)}.
\]
注意 \(p+\pi(p)=0\) 仅可能在 \(p=0\) 且 \(\pi(0)=0\) 时发生，因此每一单项式至多含一个因子 \(u_0\)。
若 \(\pi(0)=0\)，其余因子指标皆满足 \(p+\pi(p)\ge 1\)，故
\[
\Bigl|\prod_{p=0}^{\kappa-1}u_{p+\pi(p)}\Bigr|
\le
|u_0|\,B^{\kappa-1}.
\]
此类排列恰有 \((\kappa-1)!\) 个。
若 \(\pi(0)\neq 0\)，则全部指标 \(\ge 1\)，从而
\[
\Bigl|\prod_{p=0}^{\kappa-1}u_{p+\pi(p)}\Bigr|
\le
B^{\kappa},
\]
此类排列数为 \(\kappa!-(\kappa-1)!\)。
求和得第一条不等式，第二条由粗略放大系数得到。
\par
最后一条由定理 \ref{thm:xi-entropy-gap-exponential-suppression-nonzero-fingerprint} 给出 \(B\ll (\kappa_{\mathrm{ind}}-S_{\mathrm{def}})\)，并结合 \(u_0=4\pi S_{\mathrm{def}}\)（命题 \ref{prop:xi-defect-entropy-closed-form-invariant}）代入即得。证毕。
\end{proof}

\leanverified{paper\_xi\_hankel\_spike\_singular\_spectrum\_separation}
\begin{theorem}[Hankel 主块的一尖峰奇异谱分离]\label{thm:xi-hankel-spike-singular-spectrum-separation}
令 \(\mathsf H_{\kappa-1}:=(u_{p+q})_{0\le p,q\le \kappa-1}\)，并记 \(B:=\sup_{n\ge 1}|u_n|\)。
将 \(\mathsf H_{\kappa-1}\) 的奇异值按降序记为 \(\sigma_1\ge\cdots\ge\sigma_\kappa\)。
则
\[
\sigma_1(\mathsf H_{\kappa-1})\ \ge\ |u_0|-\kappa B,
\qquad
\sigma_j(\mathsf H_{\kappa-1})\ \le\ \kappa B\quad(2\le j\le \kappa).
\]
当 \(\kappa B<|u_0|\) 时，第一奇异值与其余奇异值存在显式分离。
\end{theorem}
\begin{proof}
记标准基向量 \(e_0:=(1,0,\dots,0)^{\mathsf T}\in\RR^{\kappa}\)，并分解
\[
\mathsf H_{\kappa-1}=u_0\,e_0e_0^{\mathsf T}+E,
\qquad
E_{pq}:=\begin{cases}
0,&p=q=0,\\
u_{p+q},&p+q\ge 1.
\end{cases}
\]
由定义 \(|E_{pq}|\le B\)，故
\[
\|E\|_{\mathrm{op}}
\le
\|E\|_{\mathrm{F}}
\le
\sqrt{\kappa^2-1}\,B
\le
\kappa B.
\]
秩一矩阵 \(u_0e_0e_0^{\mathsf T}\) 的奇异值为 \((|u_0|,0,\dots,0)\)。
由 Weyl 奇异值扰动不等式得
\[
\sigma_1(\mathsf H_{\kappa-1})\ge |u_0|-\|E\|_{\mathrm{op}}\ge |u_0|-\kappa B,
\qquad
\sigma_j(\mathsf H_{\kappa-1})\le \|E\|_{\mathrm{op}}\le \kappa B\ (j\ge 2).
\]
证毕。
\end{proof}

\leanverified{paper\_xi\_negative\_block\_spectrum\_spike\_and\_shallow\_bound}
\begin{corollary}[负块谱的单尖峰与浅负谱上界]\label{cor:xi-negative-block-spectrum-spike-and-shallow-bound}
令 \(G_{\kappa-1}:=\mathsf H_{\kappa-1}^{\ast}\mathsf H_{\kappa-1}\succeq 0\)。
则 \(G_{\kappa-1}\) 的特征值满足
\[
\lambda_{\max}(G_{\kappa-1})\ \ge\ (|u_0|-\kappa B)^2,
\qquad
\lambda_j(G_{\kappa-1})\ \le\ (\kappa B)^2\quad(2\le j\le \kappa).
\]
因此负块 \(-G_{\kappa-1}\) 的最深负特征值不大于 \(-(|u_0|-\kappa B)^2\)，其余负特征值皆不小于 \(-(\kappa B)^2\)。
\end{corollary}
\begin{proof}
由定义 \(\lambda_j(G_{\kappa-1})=\sigma_j(\mathsf H_{\kappa-1})^2\)，并应用定理 \ref{thm:xi-hankel-spike-singular-spectrum-separation} 即得。证毕。
\end{proof}
\begin{proposition}[浅负谱方向的秩一可消去半径]\label{prop:xi-rankone-flip-radius-shallow-negative}
在推论 \ref{cor:xi-toeplitz-scan-hankel-congruence-decomposition} 的合同正规形中，负方向由 \(-G_{\kappa-1}\) 给出。
记 \(\lambda_{\min}(G_{\kappa-1})\) 为 \(G_{\kappa-1}\) 的最小特征值，并取其单位特征向量 \(v\)。
则存在 Hermitian 秩一扰动
\(
\Delta:=\lambda_{\min}(G_{\kappa-1})\,vv^{\ast}
\)
使得相应负特征值精确平移至 \(0\)，从而负惯性至少下降 \(1\)，并且
\[
\|\Delta\|_{\mathrm{op}}=\lambda_{\min}(G_{\kappa-1})
=\sigma_\kappa(\mathsf H_{\kappa-1})^2
\le
(\kappa B)^2.
\]
若进一步使用 \(B\le 4\pi e^{-1}\,S_{\mathrm{def}}(\kappa_{\mathrm{ind}}-S_{\mathrm{def}})/\kappa_{\mathrm{ind}}\)（由命题 \ref{prop:xi-nonzero-harmonic-entropy-gap-product-sharp} 在 \(n=1\) 的特例推出），则得到由 \((S_{\mathrm{def}},\kappa_{\mathrm{ind}},\kappa)\) 参数化的显式上界。
\end{proposition}
\begin{proof}
对 Hermitian 矩阵的谱分解，沿最浅负方向的单位特征向量 \(v\) 施加秩一扰动 \(\Delta=\lambda_{\min}(G_{\kappa-1})vv^{\ast}\)，可将 \(-\lambda_{\min}(G_{\kappa-1})\) 精确推至 \(0\)，其余特征值保持不变，故负惯性至少下降 \(1\)。
范数恒等式 \(\|\Delta\|_{\mathrm{op}}=\lambda_{\min}(G_{\kappa-1})\) 由秩一算子谱半径立得。
最后由 \(\lambda_{\min}(G_{\kappa-1})=\sigma_\kappa(\mathsf H_{\kappa-1})^2\) 与定理 \ref{thm:xi-hankel-spike-singular-spectrum-separation} 的 \(\sigma_\kappa\le \kappa B\) 得到 \(\|\Delta\|_{\mathrm{op}}\le (\kappa B)^2\)。
若代入命题 \ref{prop:xi-nonzero-harmonic-entropy-gap-product-sharp}（\(n=1\)）即得显式化版本。证毕。
\end{proof}
\leanverified{paper\_xi\_rankone\_flip\_radius\_shallow\_negative}

\leanverified{paper\_xi\_entropy\_gap\_controls\_geometric\_mean\_radius}
\begin{proposition}[熵缺口控制节点模长的加权几何平均]\label{prop:xi-entropy-gap-controls-geometric-mean-radius}
在 \eqref{eq:xi-comoving-scan-fourier} 的有限缺陷口径下取 \(\Delta=1\)，并令
\[
\rho_j:=e^{-(1+\delta_j)}\in(0,e^{-1})\qquad(1\le j\le \kappa).
\]
则按重数加权的几何平均满足显式上界
\[
\Bigl(\prod_{j=1}^{\kappa}\rho_j^{m_j}\Bigr)^{1/\kappa_{\mathrm{ind}}}
\ \le\
e^{-1}\cdot\frac{\kappa_{\mathrm{ind}}-S_{\mathrm{def}}}{\kappa_{\mathrm{ind}}}.
\]
\end{proposition}
\begin{proof}
由 \(e^{\delta}\ge 1+\delta\) 得 \(e^{-\delta}\le (1+\delta)^{-1}\)，从而
\(
\rho_j=e^{-1}e^{-\delta_j}\le e^{-1}(1+\delta_j)^{-1}
\)。
对按重数展开后的 \(\kappa_{\mathrm{ind}}\) 个因子应用 AM--GM 不等式，
\[
\Bigl(\prod_{j=1}^{\kappa}(1+\delta_j)^{-m_j}\Bigr)^{1/\kappa_{\mathrm{ind}}}
\le
\frac{1}{\kappa_{\mathrm{ind}}}\sum_{j=1}^{\kappa}\frac{m_j}{1+\delta_j}
=\frac{\kappa_{\mathrm{ind}}-S_{\mathrm{def}}}{\kappa_{\mathrm{ind}}},
\]
其中最后一步使用 \eqref{eq:xi-defect-entropy-index-decomposition}。
合并即得结论。证毕。
\end{proof}

\endinput


% Spectral-gap stability of Toeplitz negative-inertia certificate under perturbations.

\begin{theorem}\label{thm:xi-toeplitz-negative-inertia-spectral-gap-stability}
\leanverified{paper\_xi\_toeplitz\_negative\_inertia\_spectral\_gap\_stability}
在有限缺陷口径下，取 \(N\) 足够大使推论 \ref{cor:xi-toeplitz-scan-hankel-congruence-decomposition} 的合同分解成立：
存在可逆矩阵 \(R_N\) 与 \(P_N\succ 0\) 使
\[
R_N^{\ast}\,\mathbf{T}_N(C)\,R_N
=
\begin{pmatrix}
P_N & 0\\
0 & -\,G_{\kappa-1}
\end{pmatrix},
\qquad
G_{\kappa-1}:=\mathsf H_{\kappa-1}^{\ast}\mathsf H_{\kappa-1}\succ 0.
\]
令 \(\widetilde{\mathbf{T}}_N:=\mathbf{T}_N(C)+E_N\) 为任意 Hermitian 扰动，并定义
\[
\Delta_N:=\bigl\|R_N^{\ast}E_N R_N\bigr\|_{\mathrm{op}},
\qquad
g_N:=\min\bigl\{\lambda_{\min}(P_N),\lambda_{\min}(G_{\kappa-1})\bigr\}.
\]
若满足谱隙条件 \(\Delta_N<g_N\)，则：
\begin{enumerate}
  \item \(\widetilde{\mathbf{T}}_N\) 的负惯性保持不变：
  \[
  \nu_-(\widetilde{\mathbf{T}}_N)=\nu_-(\mathbf{T}_N(C))=\kappa.
  \]
  \item \(\widetilde{\mathbf{T}}_N\) 在零点处保有显式谱隙：其谱与 \(0\) 的距离至少为 \(g_N-\Delta_N\)。
\end{enumerate}
\end{theorem}
\begin{proof}
由 Sylvester 惯性定理，惯性在可逆合同变换下不变，故
\[
\nu_-(\widetilde{\mathbf{T}}_N)=\nu_-\!\left(R_N^{\ast}\widetilde{\mathbf{T}}_N R_N\right).
\]
记
\[
S_N:=\begin{pmatrix}P_N&0\\0&-G_{\kappa-1}\end{pmatrix},
\qquad
F_N:=R_N^{\ast}E_N R_N,
\]
则 \(R_N^{\ast}\widetilde{\mathbf{T}}_N R_N=S_N+F_N\)，并且 \(\|F_N\|_{\mathrm{op}}=\Delta_N\)。
由于 \(P_N\succ 0\) 且 \(G_{\kappa-1}\succ 0\)，矩阵 \(S_N\) 的谱严格避开 \(0\)，且
\(
\mathrm{dist}(0,\sigma(S_N))=g_N
\)。
由 Weyl 谱扰动不等式，
每个特征值在加上 \(F_N\) 后的偏移不超过 \(\|F_N\|_{\mathrm{op}}=\Delta_N\)。
当 \(\Delta_N<g_N\) 时，正谱不穿越 \(0\)、负谱亦不穿越 \(0\)，从而惯性保持不变，并且新的谱到 \(0\) 的距离至少为 \(g_N-\Delta_N\)。
结合 Sylvester 惯性不变性即得结论。证毕。
\end{proof}

\begin{corollary}[合同尺度下的可允许扰动半径]\label{cor:xi-toeplitz-negative-inertia-allowed-perturb-radius}
\leanverified{paper\_xi\_toeplitz\_negative\_inertia\_allowed\_perturb\_radius}
在定理 \ref{thm:xi-toeplitz-negative-inertia-spectral-gap-stability} 的设定下，若进一步满足
\[
\|E_N\|_{\mathrm{op}}\ <\ \frac{g_N}{\|R_N\|_{\mathrm{op}}^{2}},
\]
则仍有 \(\nu_-(\mathbf{T}_N(C)+E_N)=\kappa\)。
并且在推论 \ref{cor:xi-toeplitz-negative-block-margin-certificate} 的设定下，负块谱隙项 \(\lambda_{\min}(G_{\kappa-1})\) 具有完全显式的下界，从而 \(g_N\) 的负向部分可由 \((w_j,z_j)\) 的分离度与有限窗口幅度 \(M\) 给出可核验证书。
\end{corollary}
\begin{proof}
由算子范数的次乘性，
\(
\|R_N^{\ast}E_N R_N\|_{\mathrm{op}}\le \|R_N\|_{\mathrm{op}}^2\|E_N\|_{\mathrm{op}}
\)。
因此若 \(\|E_N\|_{\mathrm{op}}<g_N/\|R_N\|_{\mathrm{op}}^2\)，则 \(\Delta_N<g_N\)，由定理 \ref{thm:xi-toeplitz-negative-inertia-spectral-gap-stability} 得结论。
最后一条由推论 \ref{cor:xi-toeplitz-negative-block-margin-certificate} 直接给出。证毕。
\end{proof}

\begin{corollary}[负方向的可审计保持：仅由 Hankel 负块给出的阈值]\label{cor:xi-toeplitz-negative-subspace-preservation-threshold}
\leanverified{paper\_xi\_toeplitz\_negative\_subspace\_preservation\_threshold}
在定理 \ref{thm:xi-toeplitz-negative-inertia-spectral-gap-stability} 的设定下，记
\[
\sigma_{\min}\!\bigl(\mathsf H_{\kappa-1}\bigr):=\min_{v\neq 0}\frac{\|\mathsf H_{\kappa-1}v\|_2}{\|v\|_2},
\qquad
\|R_N\|_{\mathrm{op}}:=\sup_{x\neq 0}\frac{\|R_Nx\|_2}{\|x\|_2}.
\]
若 \(E_N\) 为任意 Hermitian 扰动并满足
\[
\|E_N\|_{\mathrm{op}}\ <\ \frac{\sigma_{\min}\!\bigl(\mathsf H_{\kappa-1}\bigr)^2}{\|R_N\|_{\mathrm{op}}^{2}},
\]
则扰动后仍有
\[
\nu_-\!\bigl(\mathbf{T}_N(C)+E_N\bigr)\ \ge\ \kappa.
\]
\end{corollary}
\begin{proof}
记 \(G_{\kappa-1}=\mathsf H_{\kappa-1}^{\ast}\mathsf H_{\kappa-1}\succ 0\)，则 \(\lambda_{\min}(G_{\kappa-1})=\sigma_{\min}(\mathsf H_{\kappa-1})^2\)。
对任意 \(v\in\CC^{\kappa}\) 令 \(z:=(0\oplus v)\)，并取
\(
y:=R_N z\in\CC^{N+1}
\)。
由合同分解，
\[
y^{\ast}\mathbf{T}_N(C)y
=z^{\ast}\begin{pmatrix}P_N&0\\0&-G_{\kappa-1}\end{pmatrix}z
=-\,v^{\ast}G_{\kappa-1}v
\le -\,\lambda_{\min}(G_{\kappa-1})\|v\|_2^2.
\]
同时 \(\|y\|_2\le \|R_N\|_{\mathrm{op}}\|v\|_2\)，故
\[
y^{\ast}\bigl(\mathbf{T}_N(C)+E_N\bigr)y
\le -\,\lambda_{\min}(G_{\kappa-1})\|v\|_2^2+\|E_N\|_{\mathrm{op}}\|y\|_2^2
\le \bigl(-\lambda_{\min}(G_{\kappa-1})+\|E_N\|_{\mathrm{op}}\|R_N\|_{\mathrm{op}}^2\bigr)\|v\|_2^2<0
\]
对一切 \(v\neq 0\) 成立。
因此二次型在 \(\kappa\)-维子空间 \(\{R_N(0\oplus v):v\in\CC^\kappa\}\) 上严格负定，从而 \(\nu_-(\mathbf{T}_N(C)+E_N)\ge \kappa\)。证毕。
\end{proof}

\begin{corollary}[稳定层级的非奇异性与行列式符号]\label{cor:xi-toeplitz-determinant-sign-rigidity}
在定理 \ref{thm:xi-toeplitz-negative-inertia-spectral-gap-stability} 的设定下，Toeplitz 截断 \(\mathbf{T}_N(C)\) 必非奇异，并且其行列式符号满足
\[
\boxed{\;
\operatorname{sgn}\det\!\bigl(\mathbf{T}_N(C)\bigr)=(-1)^{\kappa}. \;}
\]
\leanverified{paper\_xi\_toeplitz\_determinant\_sign\_rigidity}
\end{corollary}
\begin{proof}
由合同分解，
\(
R_N^{\ast}\mathbf{T}_N(C)R_N=\mathrm{diag}(P_N,-G_{\kappa-1})
\)
且 \(P_N\succ 0\)、\(G_{\kappa-1}\succ 0\)，故块对角矩阵可逆，从而 \(\mathbf{T}_N(C)\) 可逆。
再由
\(
\det(R_N^{\ast}\mathbf{T}_N(C)R_N)=|\det R_N|^2\det(\mathbf{T}_N(C))
\)
知行列式符号在合同变换下保持。
而
\(
\det(\mathrm{diag}(P_N,-G_{\kappa-1}))=\det(P_N)\det(-G_{\kappa-1})
\)
的符号为 \((-1)^{\kappa}\)（因 \(G_{\kappa-1}\) 为 \(\kappa\times\kappa\) 正定）。
合并即得结论。证毕。
\end{proof}

\begin{theorem}[否证分解的秩最小性与负方向的可构造性]\label{thm:xi-toeplitz-minrank-falsifier-factorization}
\leanverified{paper\_xi\_toeplitz\_minrank\_falsifier\_factorization}
在定理 \ref{thm:xi-toeplitz-negative-inertia-spectral-gap-stability} 的设定下，令
\[
E:=\begin{pmatrix}0 & I_\kappa\end{pmatrix}\in \Mat_{\kappa\times (N+1)}(\CC)
\]
表示与合同分解的块划分相容的坐标投影。
定义
\[
Q_N:=R_N^{-*}\begin{pmatrix}P_N&0\\0&0\end{pmatrix}R_N^{-1}\succeq 0,
\qquad
W_N:=\mathsf H_{\kappa-1}\,E\,R_N^{-1}\in \Mat_{\kappa\times (N+1)}(\CC).
\]
则成立显式分解
\[
\boxed{\;\mathbf{T}_N(C)=Q_N-W_N^{\ast}W_N.\;}
\]
并且：
\begin{enumerate}
  \item \(\mathrm{rank}(W_N)=\kappa\)；
  \item \(\mathrm{Ran}(W_N^{\ast})\) 给出 \(\mathbf{T}_N(C)\) 的全部负方向：若 \(x\perp \mathrm{Ran}(W_N^{\ast})\)，则 \(x^{\ast}\mathbf{T}_N(C)x\ge 0\)；
  \item 在所有分解 \(\mathbf{T}_N(C)=Q-W^{\ast}W\)（\(Q\succeq 0\)）中，必有 \(\mathrm{rank}(W)\ge \kappa\)，从而上述分解在秩意义下最小。
\end{enumerate}
\end{theorem}
\begin{proof}
由合同分解与 \(G_{\kappa-1}=\mathsf H_{\kappa-1}^{\ast}\mathsf H_{\kappa-1}\) 得
\[
\mathbf{T}_N(C)
=R_N^{-*}\begin{pmatrix}P_N&0\\0&-\,\mathsf H_{\kappa-1}^{\ast}\mathsf H_{\kappa-1}\end{pmatrix}R_N^{-1}
=Q_N-R_N^{-*}\,E^{\ast}\mathsf H_{\kappa-1}^{\ast}\mathsf H_{\kappa-1}E\,R_N^{-1}
=Q_N-W_N^{\ast}W_N,
\]
得到分解式。
\par
第一条：因 \(R_N\) 可逆且 \(E\) 满秩，\(\mathrm{rank}(E R_N^{-1})=\kappa\)；
又 \(\mathsf H_{\kappa-1}\) 可逆，故 \(\mathrm{rank}(W_N)=\kappa\)。
\par
第二条：若 \(x\perp \mathrm{Ran}(W_N^{\ast})\)，则 \(W_Nx=0\)。
记 \(z:=R_N^{-1}x\) 并分块写 \(z=(z_+\oplus z_-)\)，则 \(W_Nx=\mathsf H_{\kappa-1}Ez=\mathsf H_{\kappa-1}z_-=0\)。
由 \(\mathsf H_{\kappa-1}\) 可逆得到 \(z_-=0\)，从而
\(
x^{\ast}\mathbf{T}_N(C)x=z^{\ast}\mathrm{diag}(P_N,-G_{\kappa-1})z=z_+^{\ast}P_N z_+\ge 0
\)。
\par
第三条：设 \(\mathbf{T}_N(C)=Q-W^{\ast}W\) 且 \(Q\succeq 0\)。
对任意 \(x\in\ker(W)\) 有 \(x^{\ast}\mathbf{T}_N(C)x=x^{\ast}Qx\ge 0\)，故任一负定子空间必包含于 \(\ker(W)^{\perp}\)。
因此
\[
\nu_-\!\bigl(\mathbf{T}_N(C)\bigr)\ \le\ \dim\ker(W)^{\perp}=\mathrm{rank}(W).
\]
而在本定理的设定下 \(\nu_-(\mathbf{T}_N(C))=\kappa\)，故 \(\mathrm{rank}(W)\ge \kappa\)，得秩最小性。证毕。
\end{proof}

\endinput


\endinput
